\chapter{Polynômes et fractions rationnelles}
\textit{Derrière les équations, les courbes et les calculs de cette classe se cachent
presque toujours des polynômes. Savoir les manipuler — déterminer leur degré, les
factoriser, les diviser, en lire le signe — est un savoir-faire central. Les fractions
rationnelles, quotients de deux polynômes, prolongent naturellement cette étude. Ce
chapitre te donne les outils que tu réutiliseras dans tout le programme.}

\section{Polynômes à une variable}

\begin{definition}[Fonction polynôme]
Une \textbf{fonction polynôme} (ou simplement \textbf{polynôme}) à une variable $x$ est
une fonction $P$ définie sur $\R$ par une expression de la forme
\[ P(x)=a_n x^n + a_{n-1}x^{n-1}+\dots+a_1 x + a_0, \]
où $n\in\N$ et où les réels $a_0,a_1,\dots,a_n$ sont appelés les \textbf{coefficients}
de $P$. Le réel $a_k$ est le coefficient du terme de degré $k$ ; $a_0$ est le
\textbf{terme constant}.
\end{definition}

\begin{definition}[Degré d'un polynôme]
Si $P$ n'est pas le polynôme nul, son \textbf{degré}, noté $\deg P$, est le plus grand
entier $n$ tel que le coefficient $a_n$ soit \emph{non nul}. Ce coefficient $a_n$ est
alors le \textbf{coefficient dominant}. Par convention, le polynôme nul n'a pas de degré.
\end{definition}

\begin{exemple}
$P(x)=2x^3-5x+7$ est un polynôme de degré $3$, de coefficient dominant $2$ et de terme
constant $7$ ; les coefficients de $x^2$ et de $x$ sont respectivement $0$ et $-5$.
Le polynôme $Q(x)=4$ est de degré $0$ (polynôme constant non nul).
\end{exemple}

\begin{propriete}[Degré d'une somme et d'un produit]
Soit $P$ et $Q$ deux polynômes non nuls. Alors :
\begin{itemize}
  \item $\deg(P\times Q)=\deg P + \deg Q$ ;
  \item $\deg(P+Q)\le \max(\deg P,\deg Q)$, avec égalité dès que $\deg P\neq \deg Q$.
\end{itemize}
\end{propriete}

\begin{remarque}
Le degré d'une somme peut \emph{baisser} si les termes dominants se compensent. Par
exemple, $(x^2+x)+(-x^2+1)=x+1$ : deux polynômes de degré $2$ dont la somme est de
degré $1$.
\end{remarque}

\begin{theoreme}[Égalité de deux polynômes]
Deux polynômes sont \textbf{égaux} si et seulement s'ils ont \emph{le même degré} et
\emph{les mêmes coefficients}, terme à terme. Autrement dit, si
\[ a_n x^n+\dots+a_0 = b_n x^n+\dots+b_0 \quad\text{pour tout }x\in\R, \]
alors $a_k=b_k$ pour tout $k$. En particulier, un polynôme est nul (pour tout $x$) si et
seulement si \emph{tous} ses coefficients sont nuls.
\end{theoreme}

\begin{aretenir}[Identification des coefficients]
Cette propriété fonde la méthode d'\textbf{identification} : pour déterminer des
coefficients inconnus, on développe, puis on égale terme à terme les coefficients de
même degré des deux membres. On obtient alors un système d'équations.
\end{aretenir}

\begin{exemple}
Cherchons $a,b,c$ tels que $x^2+1=a(x-1)^2+b(x-1)+c$ pour tout $x$. En développant le
membre de droite : $a x^2+(-2a+b)x+(a-b+c)$. L'identification donne
\[ a=1,\qquad -2a+b=0,\qquad a-b+c=1, \]
d'où $a=1$, $b=2$, $c=2$. Ainsi $x^2+1=(x-1)^2+2(x-1)+2$.
\end{exemple}

\section{Racines et factorisation}

\begin{definition}[Racine d'un polynôme]
On appelle \textbf{racine} (ou \textbf{zéro}) du polynôme $P$ tout réel $a$ tel que
$P(a)=0$.
\end{definition}

\begin{theoreme}[Racine et factorisation]
Soit $P$ un polynôme et $a$ un réel. Alors :
\[ a \text{ est racine de } P \iff (x-a) \text{ divise } P, \]
c'est-à-dire qu'il existe un polynôme $Q$ tel que $P(x)=(x-a)\,Q(x)$ pour tout $x$.
\end{theoreme}

\begin{remarque}
Démonstration du sens direct. Si $a$ est racine, écrivons $P(x)=P(x)-P(a)$ (car
$P(a)=0$). Chaque terme $a_k(x^k-a^k)$ se factorise par $(x-a)$, puisque
$x^k-a^k=(x-a)(x^{k-1}+ax^{k-2}+\dots+a^{k-1})$. En mettant $(x-a)$ en facteur dans
toute la somme, on obtient $P(x)=(x-a)Q(x)$. Réciproquement, si $P(x)=(x-a)Q(x)$, alors
$P(a)=(a-a)Q(a)=0$, donc $a$ est racine.
\end{remarque}

\begin{propriete}[Polynôme factorisé et nombre de racines]
Si $a_1,a_2,\dots,a_p$ sont des racines \emph{distinctes} de $P$, alors $P$ se factorise
par le produit $(x-a_1)(x-a_2)\cdots(x-a_p)$. En conséquence, un polynôme non nul de
degré $n$ admet \textbf{au plus $n$ racines}.
\end{propriete}

\begin{exemple}
$P(x)=x^3-3x^2+2x$. On remarque $P(0)=0$, donc $x$ se met en facteur :
$P(x)=x(x^2-3x+2)$. Or $1$ est racine de $x^2-3x+2$, donc
$x^2-3x+2=(x-1)(x-2)$. Finalement $P(x)=x(x-1)(x-2)$ : les racines sont $0$, $1$ et $2$.
\end{exemple}

\begin{aretenir}[Trouver une racine évidente]
Pour amorcer une factorisation, on teste les valeurs simples : $0$, $1$, $-1$, $2$,
$-2$\dots\ Si $P(a)=0$, alors $(x-a)$ est en facteur ; on obtient le quotient par
division euclidienne ou par identification.
\end{aretenir}

\section{Division des polynômes}

\subsection{Division euclidienne}

\begin{theoreme}[Division euclidienne]
Soit $A$ et $B$ deux polynômes avec $B\neq 0$. Il existe un \emph{unique} couple de
polynômes $(Q,R)$ tel que
\[ A = B\,Q + R \qquad\text{avec}\qquad R=0 \ \text{ ou } \ \deg R < \deg B. \]
$Q$ est le \textbf{quotient} et $R$ le \textbf{reste} de la division euclidienne de $A$
par $B$. On dit que $B$ \textbf{divise} $A$ lorsque $R=0$.
\end{theoreme}

\begin{aretenir}[Poser la division]
On dispose la division comme pour les entiers, en ordonnant $A$ et $B$ suivant les
\textbf{puissances décroissantes}. À chaque étape, on divise le terme dominant du reste
courant par le terme dominant de $B$, on multiplie, on soustrait, et on recommence tant
que le degré du reste reste $\ge \deg B$.
\end{aretenir}

\begin{exemple}
Divisons $A(x)=2x^3-3x^2+4x-1$ par $B(x)=x^2+1$.
\begin{itemize}
  \item $2x^3 \div x^2 = 2x$ ; $\ 2x(x^2+1)=2x^3+2x$ ; reste $-3x^2+2x-1$.
  \item $-3x^2 \div x^2 = -3$ ; $\ -3(x^2+1)=-3x^2-3$ ; reste $2x+2$.
\end{itemize}
Le degré de $2x+2$ étant $<2$, on s'arrête. Ainsi
\[ 2x^3-3x^2+4x-1=(x^2+1)(2x-3)+(2x+2), \]
avec quotient $Q(x)=2x-3$ et reste $R(x)=2x+2$.
\end{exemple}

\subsection{Division suivant les puissances décroissantes}

\begin{remarque}
La division précédente s'appelle aussi \textbf{division suivant les puissances
décroissantes} : on classe les termes du plus haut degré au plus bas, et on s'arrête dès
que le degré du reste devient strictement inférieur à celui du diviseur. C'est la forme
employée pour factoriser un polynôme et pour décomposer une fraction rationnelle.
\end{remarque}

\begin{exemple}
On sait que $1$ est racine de $A(x)=x^3-2x^2-x+2$. Divisons par $(x-1)$ :
\[ x^3-2x^2-x+2=(x-1)(x^2-x-2). \]
Puis $x^2-x-2=(x-2)(x+1)$. D'où la factorisation complète
$A(x)=(x-1)(x-2)(x+1)$, de racines $1$, $2$, $-1$.
\end{exemple}

\section{Identités remarquables et factorisations}

\begin{propriete}[Identités remarquables]
Pour tous réels $a$ et $b$ :
\begin{align*}
(a+b)^2 &= a^2+2ab+b^2, & (a-b)^2 &= a^2-2ab+b^2,\\
(a+b)(a-b) &= a^2-b^2, & (a+b)^3 &= a^3+3a^2b+3ab^2+b^3,\\
(a-b)^3 &= a^3-3a^2b+3ab^2-b^3. &&
\end{align*}
Et pour les sommes/différences de cubes :
\[ a^3-b^3=(a-b)(a^2+ab+b^2),\qquad a^3+b^3=(a+b)(a^2-ab+b^2). \]
\end{propriete}

\begin{exemple}
$x^3-8=x^3-2^3=(x-2)(x^2+2x+4)$. Le facteur $x^2+2x+4$ n'a pas de racine réelle (son
discriminant vaut $4-16=-12<0$), donc cette factorisation est la plus poussée dans $\R$.
\end{exemple}

\begin{aretenir}[Stratégie de factorisation]
\begin{itemize}
  \item mettre d'abord en \textbf{facteur commun} ce qui est commun à tous les termes ;
  \item reconnaître une \textbf{identité remarquable} ;
  \item chercher une \textbf{racine évidente}, puis factoriser par $(x-a)$ via une division ;
  \item pour un trinôme $ax^2+bx+c$, utiliser les racines (voir section suivante).
\end{itemize}
\end{aretenir}

\section{Signe d'un polynôme}

\begin{propriete}[Signe du trinôme du second degré]
Soit $P(x)=ax^2+bx+c$ avec $a\neq 0$, de discriminant $\Delta=b^2-4ac$.
\begin{itemize}
  \item Si $\Delta>0$ : $P$ a deux racines $x_1<x_2$ ; $P(x)$ est \emph{du signe de $a$}
        à l'extérieur de $[x_1\,;\,x_2]$ et \emph{du signe de $-a$} entre les racines.
  \item Si $\Delta=0$ : $P$ a une racine double $x_0=-\dfrac{b}{2a}$ ; $P(x)$ est du
        signe de $a$ partout, et s'annule en $x_0$.
  \item Si $\Delta<0$ : $P$ n'a pas de racine réelle ; $P(x)$ est du signe de $a$ pour
        tout $x$.
\end{itemize}
On résume cela par la formule : « \textbf{le trinôme est du signe de $a$ sauf entre les
racines} ».
\end{propriete}

\begin{illus}
\begin{tikzpicture}[scale=0.95]
  \draw[->,onbline] (-1.4,0)--(4.4,0) node[right,onbmuted]{$x$};
  \draw[->,onbline] (0,-1.6)--(0,3.2) node[above,onbmuted]{$y$};
  \draw[thick,onbo,smooth,samples=80,domain=-0.8:3.8] plot (\x,{(\x-1)*(\x-3)});
  \filldraw[onbgreen] (1,0) circle (1.7pt) node[below right,onbink]{$x_1$};
  \filldraw[onbgreen] (3,0) circle (1.7pt) node[below right,onbink]{$x_2$};
  \node[onbmuted] at (2,-1.25) {$P<0$};
  \node[onbmuted] at (-0.6,2.4) {$P>0$};
  \node[onbmuted] at (3.7,2.4) {$P>0$};
\end{tikzpicture}
\fig{Figure — pour $a>0$ (ici $P(x)=(x-1)(x-3)$) : $P$ est négatif entre les racines, positif à l'extérieur.}
\end{illus}

\begin{propriete}[Signe d'un polynôme factorisé]
Lorsqu'un polynôme est écrit sous forme de \textbf{produit de facteurs du premier
degré} (et éventuellement de trinômes sans racine), on dresse un \textbf{tableau de
signes} : chaque facteur $(x-a)$ est négatif pour $x<a$ et positif pour $x>a$ ; le signe
du produit s'obtient par la règle des signes.
\end{propriete}

\begin{illus}
\renewcommand{\arraystretch}{1.35}
\begin{tabular}{@{}c|ccccccc@{}}
\toprule
$x$        & $-\infty$ && $-1$ && $2$ && $+\infty$ \\
\midrule
$x-2$      &           & $-$  & $-$ & $-$ & $0$ & $+$ & \\
$x+1$      &           & $-$  & $0$ & $+$ & $+$ & $+$ & \\
\midrule
$(x-2)(x+1)$ &         & $+$  & $0$ & $-$ & $0$ & $+$ & \\
\bottomrule
\end{tabular}
\fig{Figure — tableau de signes de $P(x)=(x-2)(x+1)=x^2-x-2$. Le produit est négatif sur $]-1\,;\,2[$.}
\end{illus}

\begin{exemple}
Résolvons $x^2-x-2\le 0$. D'après le tableau ci-dessus, $P(x)=(x-2)(x+1)$ est négatif ou
nul exactement sur $[-1\,;\,2]$. L'ensemble des solutions est donc $S=[-1\,;\,2]$.
\end{exemple}

\section{Fractions rationnelles}

\begin{definition}[Fraction rationnelle]
On appelle \textbf{fraction rationnelle} tout quotient $F(x)=\dfrac{P(x)}{Q(x)}$ de deux
polynômes $P$ et $Q$, avec $Q$ non nul. L'\textbf{ensemble de définition} de $F$ est
l'ensemble des réels qui n'annulent pas le dénominateur :
\[ D_F=\{\,x\in\R \ \mid\ Q(x)\neq 0\,\}. \]
\end{definition}

\begin{exemple}
Pour $F(x)=\dfrac{x+3}{x^2-4}$, le dénominateur s'annule en $x=2$ et $x=-2$. Donc
\[ D_F=\,]-\infty\,;\,-2[\,\cup\,]-2\,;\,2[\,\cup\,]2\,;\,+\infty[. \]
\end{exemple}

\begin{aretenir}[Simplifier une fraction rationnelle]
On factorise le numérateur et le dénominateur, puis on simplifie le ou les facteurs
communs. \textbf{Attention} : la simplification ne change pas $D_F$ — les valeurs
interdites du dénominateur \emph{initial} restent exclues.
\end{aretenir}

\begin{exemple}
$F(x)=\dfrac{x^2-1}{x^2-x}=\dfrac{(x-1)(x+1)}{x(x-1)}$. On a $D_F=\R\setminus\{0\,;\,1\}$.
Après simplification par $(x-1)$ : $F(x)=\dfrac{x+1}{x}$ pour tout $x\in D_F$. Le réel
$1$ reste exclu même s'il ne fait plus apparaître de division par zéro après simplification.
\end{exemple}

\begin{propriete}[Signe d'une fraction rationnelle]
Le signe de $F(x)=\dfrac{P(x)}{Q(x)}$ s'étudie dans un \textbf{tableau de signes}
regroupant les signes de $P$ et de $Q$, en appliquant la règle des signes du quotient.
Aux valeurs annulant $Q$, $F$ \textbf{n'est pas définie} : on y place une double barre.
\end{propriete}

\begin{exemple}
Étudions le signe de $F(x)=\dfrac{x-1}{x+2}$. Le numérateur s'annule en $1$, le
dénominateur en $-2$ (valeur interdite). On lit : $F(x)>0$ sur
$]-\infty\,;\,-2[\cup\,]1\,;\,+\infty[$, $F(x)<0$ sur $]-2\,;\,1[$, et $F(1)=0$.
\end{exemple}

\begin{remarque}
Une fraction rationnelle peut parfois s'écrire comme \textbf{somme d'un polynôme et
d'une fraction plus simple}. On y parvient par une division euclidienne du numérateur par
le dénominateur. Par exemple :
$\dfrac{x^2+1}{x-1}=\dfrac{(x-1)(x+1)+2}{x-1}=x+1+\dfrac{2}{x-1}$.
Cette décomposition facilite l'étude du comportement de la fraction.
\end{remarque}

% ════════════════════════════════════════════════════════════
\section{L'essentiel à retenir}

\begin{synthese}
\textbf{Polynôme.} $P(x)=a_n x^n+\dots+a_0$ ; degré = plus haut exposant à coefficient
non nul. \textbf{Égalité} : même degré et mêmes coefficients (d'où l'\emph{identification}).

\smallskip
\textbf{Racine et facteur.} $a$ racine de $P \iff P(a)=0 \iff (x-a)$ divise $P$.
Un polynôme de degré $n$ a au plus $n$ racines.

\smallskip
\textbf{Division euclidienne.} $A=BQ+R$ avec $R=0$ ou $\deg R<\deg B$ ; $B$ divise $A$ si $R=0$.

\smallskip
\textbf{Identités.} $a^2-b^2=(a-b)(a+b)$ ; $a^3-b^3=(a-b)(a^2+ab+b^2)$ ;
$a^3+b^3=(a+b)(a^2-ab+b^2)$.

\smallskip
\textbf{Signe du trinôme} $ax^2+bx+c$ : du signe de $a$ \emph{sauf entre les racines}
(si $\Delta>0$) ; toujours du signe de $a$ si $\Delta\le 0$.

\smallskip
\textbf{Fraction rationnelle} $\dfrac{P}{Q}$ : $D_F$ exclut les zéros de $Q$ ;
simplifier sans changer $D_F$ ; signe via un tableau (double barre aux valeurs interdites).
\end{synthese}

% ════════════════════════════════════════════════════════════
\section{Méthodes \& savoir-faire}

\methode{Déterminer des coefficients par identification}
Développer le membre où figurent les inconnues, puis égaler terme à terme les
coefficients de même degré des deux membres ; résoudre le système obtenu.
\begin{exemple}
Trouvons $a,b$ tels que $\dfrac{3x-1}{x^2}=\dfrac{a}{x}+\dfrac{b}{x^2}$. En réduisant au
même dénominateur : $\dfrac{ax+b}{x^2}$. L'identification donne $a=3$ et $b=-1$, soit
$\dfrac{3x-1}{x^2}=\dfrac{3}{x}-\dfrac{1}{x^2}$.
\end{exemple}

\methode{Factoriser un polynôme à partir d'une racine évidente}
Tester $0,1,-1,2,-2$\dots\ Dès qu'une valeur $a$ annule $P$, effectuer la division
euclidienne de $P$ par $(x-a)$, puis factoriser le quotient.
\begin{exemple}
$P(x)=x^3+2x^2-x-2$. On a $P(1)=1+2-1-2=0$, donc $(x-1)$ divise $P$. La division donne
$P(x)=(x-1)(x^2+3x+2)=(x-1)(x+1)(x+2)$.
\end{exemple}

\methode{Effectuer une division euclidienne de polynômes}
Ordonner suivant les puissances décroissantes, puis diviser le terme dominant du reste
par celui du diviseur, multiplier, soustraire, recommencer tant que $\deg(\text{reste})\ge\deg B$.
\begin{exemple}
$A(x)=x^3+0x^2+0x+1$ par $B(x)=x+1$. On obtient $x^3+1=(x+1)(x^2-x+1)$, reste nul :
$x+1$ divise $x^3+1$ (c'est l'identité $a^3+b^3$ avec $b=1$).
\end{exemple}

\methode{Étudier le signe d'un trinôme}
Calculer $\Delta=b^2-4ac$. Si $\Delta\ge 0$, trouver les racines ; appliquer la règle
« du signe de $a$ sauf entre les racines ». Si $\Delta<0$, le trinôme garde le signe de $a$.
\begin{exemple}
$P(x)=-2x^2+3x+2$. $\Delta=9+16=25>0$, racines
$x=\dfrac{-3\pm5}{-4}$, soit $x_1=-\dfrac12$ et $x_2=2$. Comme $a=-2<0$, $P>0$ entre
$-\frac12$ et $2$, et $P<0$ à l'extérieur.
\end{exemple}

\methode{Déterminer l'ensemble de définition d'une fraction rationnelle}
Factoriser le dénominateur et résoudre $Q(x)=0$ ; exclure ces valeurs de $\R$.
\begin{exemple}
$F(x)=\dfrac{2x}{x^2-3x}=\dfrac{2x}{x(x-3)}$. Le dénominateur s'annule en $0$ et $3$,
donc $D_F=\R\setminus\{0\,;\,3\}$. (Après simplification, $F(x)=\dfrac{2}{x-3}$ sur $D_F$.)
\end{exemple}

\methode{Étudier le signe d'une fraction rationnelle par un tableau}
Factoriser numérateur et dénominateur ; placer toutes les valeurs critiques dans un
tableau ; lire le signe de chaque facteur puis appliquer la règle des signes du quotient,
avec une double barre aux valeurs interdites.
\begin{exemple}
$F(x)=\dfrac{x+1}{x-2}$. Valeurs critiques : $-1$ (zéro) et $2$ (interdite). Le tableau
donne $F>0$ sur $]-\infty\,;\,-1[\cup\,]2\,;\,+\infty[$ et $F<0$ sur $]-1\,;\,2[$, avec
$F(-1)=0$ et $F$ non définie en $2$.
\end{exemple}
